\section{Experiments, predictive realization and admissible resources}\label{sec:experiments}
\subsection{Primitive data and predictive quotient}
All quotient assertions are made at a fixed experiment interface phase. In a time-inhomogeneous experiment the phase belongs to that interface; an autonomous controller must still carry its own phase through its charged labels.

A prepared causal experiment consists of standard Borel report spaces $Y_a$, a finite nonempty action set $A$, positive causal probability kernels for successive reports and protected terminal tests, and a preparation kernel. A legal strategy selects the next action from the observed history and independent private randomness. A future test is a bounded function of a legally executable finite continuation, including its prescribed interventions and terminal readout. Two prepared histories are equivalent when all such future test expectations coincide. Convex preparation means an actual random mixture of preparations, not a formal interpolation of an unimplemented state.

We impose the following finite realization hypothesis. A finite collection of executable tests and the constant test spans a separating affine space of dimension $d+1$. Their expectation vector $x(h)$ is Borel. Its attained image is Borel, with a Borel realization section. The expectations of all executable future tests factor through this vector. The unnormalized joint expectation of any coordinate test and the next report lies in the same affine span. A compact convex set $D\subset\R^d$, with nonempty interior in these coordinates, consists of realizable mixed preparations; all these affine identities hold on $D$. Norms and volumes below refer to this affine test chart, not to an ambient physical parameterization.

Consequently, for each $a$ and phase $t$ there are a reference probability measure $\lambda_{t,a}$ and Borel affine functions
\begin{equation}\label{eq:raw}
 \ell_{t,a}(y,u)=\ell^0_{t,a}(y)+c_{t,a}(y)u,\qquad
 U_{t,a}(y,u)=b_{t,a}(y)+A_{t,a}(y)u
\end{equation}
such that the report law is $\ell_{t,a}(y,u)\lambda_{t,a}(dy)$ and the normalized predictive update is
\begin{equation}\label{eq:posterior}
 F_{t,a}(y,u)=U_{t,a}(y,u)/\ell_{t,a}(y,u).
\end{equation}
Here $c_{t,a}(y)$ is a row linear functional and $c_{t,a}(y)v$ is scalar. For each $(t,a)$, the identities hold for all $u\in D$ off \emph{one} reference-null report set. When $\ell=0$, positivity and bounded coordinate tests force $U=0$ at the relevant state; the identity-level formulas use the closed perspective convention. The sharp theorem below uses positive common support. This hypothesis is verified from raw kernels in Sections~\ref{sec:classical} and~\ref{sec:quantum}; finite physical dimension alone would not prove it.

\begin{lemma}[Borel predictive realization and minimality]\label{lem:quotient}
Under the finite realization hypothesis, the attained test image, with its Borel structure, realizes the predictive quotient. Every other Borel statistic sufficient for all executable future tests determines this image by a Borel map on its realized image, provided it has a Borel realization section. The report and update kernels descend jointly in the finite action coordinate.
\end{lemma}
\begin{proof}
Equality of test coordinates implies equality of all future tests by separation and the factorization hypothesis; the converse follows because the coordinates themselves are executable. The Borel section identifies the quotient with the attained image, not with all of $D$. For another statistic $z$, compose its realization section with $h\mapsto x(h)$; sufficiency makes this independent of the representative. This gives the asserted Borel factorization. Formula~\eqref{eq:raw}, the common null set, and the finite action set give jointly Borel report and normalized update kernels. Iterated kernel integration reproduces every finite strategy and terminal test, proving realization. No unproved assertion that every measurable quotient is standard Borel is used.
\end{proof}

The larger $D$ is a domain of realizable mixtures used for uniform certificates. It is not asserted to be the actually acquired support. For an incoming $u$ define the acquired row law
\begin{equation}\label{eq:row}
 \mu_{t,a,u}=(F_{t,a}(\cdot,u))_\#
       (\ell_{t,a}(\cdot,u)\lambda_{t,a}).
\end{equation}
Under a full-history strategy $\pi$, the actual predictive vector $X_t$ and its law $\nu_t^\pi$ are obtained by iterating the chosen kernels. Under a finite controller, the relevant row mixture instead uses that controller's actual conditional barycenters. These two laws need not agree.

\subsection{The task and the two memory architectures}
The terminal audit is a bounded executable vector $Z$ satisfying
\begin{equation}\label{eq:task}
 \E[Z\mid\text{predictive state }u]=u,\qquad
 \E[\norm Z^2\mid u]=c(u),\qquad
 g(u)=c(u)-\norm u^2,
\end{equation}
where $c$ is affine. The response $r\in D$ is committed before the protected audit, and the loss is $\norm{Z-r}^2$. We also allow an audit in a larger Euclidean space with mean $i(u)$, where $i(u)=z_0+Ju$ is an affine isometry of the predictive chart. For loss $\norm{Z-i(r)}^2$, put $c(u)=\E[\norm{Z-z_0}^2\mid u]$. Direct expansion gives the same conditional loss $c(u)-2\langle u,r\rangle+\norm r^2$ and the same $g$. Thus audit noise orthogonal to the mean image is retained in the affine second moment, not discarded. This identifies the ordinary categorical or informationally complete POVM Brier task exactly with the displayed algebra; an alternative affine chart only changes fixed norm constants. All risks use this one physical task metric.

For fixed $n$, a checkpoint strategy may use the entire history to choose its $n$ acquisitions, then replace that history by at most $m$ labels before the audit. Its acquisition workspace is deliberately unrestricted and is not identified with the checkpoint label budget. Write its optimum as $\calR_{\rm cp}(n,m)$.

An autonomous controller has one finite set $S$ of at most $M$ labels, one initial label, a halt status and action kernel at every nonhalt label, a Borel next-label kernel depending only on the current label, chosen action and fresh report, and a terminal response attached to each halt label. All its phase and control information is in $S$. Fresh random coins are discarded after the update; a retained random seed would count as state. There is no external time, external last-report signal, re-readable history, persistent real register or free interim response register. The machine must halt after exactly $n$ reports almost surely. The underlying experiment supplies further reports if requested, so the horizon is not signaled by ending the physical input. Its optimum is $\calR_{\rm aut}(n,M)$; an empty class has value $+\infty$.

Known experiment constants and programmed conditional centroids are read-only program data, not free bit complexity. This definition concerns hard-state existence. Section~\ref{sec:resources} specifies additional precision, program, workspace and primitive-cost certificates for finite-bit execution. Preparation, each raw report, and the protected audit are separate physical calls: the core count is $N=n+2$.

The full-history optimum is given by
\begin{equation}\label{eq:bellman}
 V_n=g,\qquad V_{t-1}(u)=\min_{a\in A}\int V_t(w)\mu_{t,a,u}(dw).
\end{equation}
All full-history randomized strategies for the same preparation, action family and audit are admitted. The positivity and continuity conditions below ensure Borel minimizing selectors. A fixed strategy has an affine physical risk before optimization over its terminal response, so $V_t$ is concave. Thus~\eqref{eq:bellman} is a control optimum, not a same-controller substitution of approximate beliefs inside a cost.
